\documentclass[output=paper,colorlinks,citecolor=brown]{langscibook}
\ChapterDOI{10.5281/zenodo.22880226}
\title{The Teop people and their language}
\author{Ulrike Mosel}
\abstract{\noabstract}
\IfFileExists{../localcommands.tex}{
   \addbibresource{../localbibliography.bib}
   \usepackage{langsci-optional}
\usepackage{langsci-gb4e}
\usepackage{langsci-lgr}

\usepackage{listings}
\lstset{basicstyle=\ttfamily,tabsize=2,breaklines=true}

% \usepackage{noto}
% \usepackage{graphicx}
% \usepackage{hyperref}
% \usepackage{booktabs}
\usepackage{multicol}
\usepackage{array}
\usepackage{vowel}
% \usepackage{tabularx}
% \newcolumntype{Q}{>{\raggedright\arraybackslash}X}
% \usepackage{longtable}
\usepackage{rotating}
\usepackage{lscape}
\usepackage{wasysym}
\usepackage{paracol}
\usepackage{parallel}
\usepackage{listofitems}
\usepackage{langsci-branding}
% \usepackage[inline]{showlabels}
% \renewcommand{\showlabelfont}{\small\color{blue}}
\usepackage{xltabular}
\usepackage{multirow}
\usepackage{hanging}
\usepackage{graphicx}
\usepackage{subfig}
\usepackage{caption}
\usepackage{xurl}

   % \newcommand{\addtext}[2]{
%     \chapter{#2}\label{text:#1}
%     \section*{Metadata}
%     \input{chapters/#1/meta_#1.tex}
%     \section*{Summary \& background}
%     \input{chapters/#1/intro_#1.tex}
%     \section*{Parallel Text}
%     \input{chapters/#1/translation_#1.tex}
%     \section*{Glossed text}
%     \input{chapters/#1/glosses_#1.tex}
% }

\newcommand{\til}{$\sim$}

\hypersetup{hypertexnames=false}%fixes the wrong hyperlinks after \setcounter{chapter}{0} in main.tex

\newcommand{\textref}[1]{Text~\ref{#1}}%xref command for texts
\renewcommand{\sectref}[1]{Section~\ref{#1}}%xref command sections

%\examplesitalics%morpheme line of example appears in italics
\renewcommand{\exfont}{\itshape}%text line of examples appears in italics


\DeclareGraphicsExtensions{.jpg,.png,.pdf}

% \RenewDocumentCommand{\includepaper}{m}{%
%     \bgroup
%       \renewcommand{\newlineCover}{\\}
%       \renewcommand{\documentclass}[2][]{}%
%       \renewcommand{\usepackage}[2][]{}%
%       \renewenvironment{document}{\begingroup}{\endgroup}
%       \includepaper@body
%       % \begin{refsection}
%         \renewcommand{\lsCollectionPaperCitationText}{\fullciteFooter{#1footer}}
%         \include{#1}%
%         \if@partsw\relax\else% This switch controls whether the included chapter is in the range of \includeonly. It's from source2e.
%           \addtocounter{page}{-1}
%           \edef\lsCollectionPaperLastPage{\thepage}	% \lsCollectionPaperFirstPage is defined in \includepaper@body
%           \addtocounter{page}{1}
%           \lsCollectionMetadataToBibliography{#1}
%         \fi% If the paper is not within \includeonly, don't do anything.
%       % \end{refsection}
%     \egroup
% } %end NewDocumentCommand

   %% hyphenation points for line breaks
%% Normally, automatic hyphenation in LaTeX is very good
%% If a word is mis-hyphenated, add it to this file
%%
%% add information to TeX file before \begin{document} with:
%% %% hyphenation points for line breaks
%% Normally, automatic hyphenation in LaTeX is very good
%% If a word is mis-hyphenated, add it to this file
%%
%% add information to TeX file before \begin{document} with:
%% %% hyphenation points for line breaks
%% Normally, automatic hyphenation in LaTeX is very good
%% If a word is mis-hyphenated, add it to this file
%%
%% add information to TeX file before \begin{document} with:
%% \include{localhyphenation}
\hyphenation{
    par-a-digm
    Bougain-ville
    Mela-nes-ian
    Meso-Mela-nes-ian
}

\hyphenation{
    par-a-digm
    Bougain-ville
    Mela-nes-ian
    Meso-Mela-nes-ian
}

\hyphenation{
    par-a-digm
    Bougain-ville
    Mela-nes-ian
    Meso-Mela-nes-ian
}

   %\boolfalse{bookcompile}
   %\togglepaper[23]%%chapternumber
   }
\begin{document}
\maketitle\label{chap:1}
The Teop people live in the Autonomous Region of Bougainville Island of Papua New Guinea, in the Local Level Government area Tinputz of the North Bougainville District. On the map in \figref{fig:map} this area includes the Teop Harbour which lies south of the village Tinputz.

\begin{figure}[ht]
    \centering
    \includegraphics[width=.88\linewidth]{figures/map.png}
    \caption{Vicinity of Bougainville Island (from \url{https://commons.wikimedia.org/wiki/File:Un-bougainville.png})}
    \label{fig:map}
\end{figure}

Teop is an Oceanic Austronesian language that belongs to the Nehan-North Bougainville network of the North-West Solomonic Group of the Meso-Mela\-nesian North-West Solomonic Group (\cite[217]{Ross:1988ai}).\footnote{Extensive resources in and about the Teop language are archived in the Open Language Archives Community (OLAC), see \url{http://language-archives.org/} and OLAC resources in and about the Teop language: \url{http://www.language-archives.org/language/tio}.} \figref{fig:languagemap} shows the languages spoken in the Autonomous Region of Bougainville. The current number of speakers of the Teop language is unknown. According to \textcite{Grimes:2000bz}, quoted by \textcite[37]{Tryon:2015sq} the number of speakers was 5,000 in 1991. During our documentation project (1994--2014) Teop people estimated that the number of Teop speakers was about 6,000. 

\begin{figure}[ht]
    \centering
    \includegraphics[width=.8\linewidth]{figures/CartoGIS-04-006_T1.jpg}
    \caption{The languages of the Autonomous Region of Bougainville, Papua New Guinea (from ANU CartoGIS)}
    \label{fig:languagemap}
\end{figure}

Besides the official languages English and Tok Pisin, the Autonomous Region of Bougainville counts 17 Austronesian and eight Papuan languages.\footnote{See \textcite[37, 40]{Tryon:2015sq} and \textcite[2--4]{Smith-Dennis:2020mu}, and the references given there.} According to the Bougainville Human Development Report more than 80\% of the population of North Bougainville were literate (\cite{Chalmers:2009qz}). The report does not distinguish  between literacy in English, Tok Pisin or vernacular languages in its statistics.
\end{document}